\section{A common master and birational transport}
\label{sec:global-transport}

The obstruction in Sections~\ref{sec:incomplete-gamma}--\ref{sec:two-wall}
comes from comparing two local receivers at an intermediate chamber.  We now
avoid that comparison.  A single equivariant cobordism carries both endpoint
quantum Kirwan maps, and virtual localization compares their point rows before
either endpoint Novikov completion is taken.

\subsection{The common point row}

Let \(W\) be a smooth projective variety with a \(\Gm\)-action and let
\(Y_-\) and \(Y_+\) be smooth free GIT quotients for two extreme
linearizations.  We assume the stable-equals-semistable and properness
hypotheses under which Woodward's localized gauged theory is defined.  In
particular, degree by degree, the polarization master stack is a proper
Deligne--Mumford stack with its stated perfect obstruction theory.  We work
coefficientwise at sufficiently large area.  Write \(\kappa_\pm\) for
the corresponding quantum Kirwan maps and \(\tau_{Y_\pm,-}\) for Woodward's
localized graph potentials.  The
derivatives
\[
 A_\pm=D\tau_{Y_\pm,-}\,D\kappa_\pm                       \tag{8.1}
\]
are the endpoint localized gauged maps; this is the localized adiabatic
identity \cite[Equation~(68)]{WoodwardQKIII}.

Suppose that a class \(a_p\in H^*_{\Gm}(W)\) has the following restrictions:
\[
 \kappa_\pm(a_p)=[p_\pm],\qquad
 a_p|_F=0                                                     \tag{8.2}
\]
for every fixed component \(F\) occurring at an intermediate polarization.
Here \(p_-\) and \(p_+\) lie in the common birational open set.

\begin{proposition}[Support collapse]
\label{prop:support-collapse}
At every fixed equivariant degree, the point-marked virtual Kalkman formula
has no intermediate-polarization contribution.  Hence its two endpoint
coefficients agree in the extended Novikov coefficient-function space.
After Rees homogenization and in any common realization of the two endpoint
series, this reads
\[
 \rrow_{Y_-,p_-}A_- =\rrow_{Y_+,p_+}A_+.                    \tag{8.3}
\]
The identity is coefficientwise in all remaining Novikov and bulk
variables; the proposition does not construct the common realization.
\end{proposition}

\begin{proof}
Retain rotation of the parametrized graph curve, with equivariant parameter
\(\hbar\).  At each ordinary input marking insert the normalized equivariant
point class of \(0\in\mathbf P^1\).  Rotation localization then puts every
input on the zero-side bubble tree.  Next choose equivariant divisor classes
\(D_1,\ldots,D_s\) spanning the full equivariant numerical dual of the
allowed total degrees; their endpoint Kirwan restrictions in particular span
the endpoint numerical divisor spaces.  Insert Woodward's Liouville classes
with parameters \(t_j\).  On a graph fixed locus their contribution is
\[
 \exp\!\left(\sum_jt_jD_j\right)
 \prod_j x_j^{D_j\cdot \widetilde d_+},
 \qquad x_j=\exp(t_j\hbar),                                  \tag{8.4}
\]
where \(\widetilde d_+\) is the total infinity-side equivariant degree: the
bubble degree together with the affine cocharacter degree of the gauged
principal component \cite[Corollary~9.10(c)]{WoodwardQKIII}.  Extracting the
trivial character in all \(x_j\) forces \(\widetilde d_+=0\).  The infinity
factor is therefore the unstable identity: positivity on the effective
infinity-side semigroup leaves no nonconstant degree-zero component.  The
zero factor is exactly the localized graph potential in (8.1).  Using the
full equivariant divisor lattice is important: endpoint divisor spaces alone
need not separate an affine cocharacter degree.

For each fixed degree and sufficiently large area, every nonendpoint
polarization-fixed graph has its principal component in a fixed quotient.
Moreover every marking and node lands in the corresponding fixed locus
\cite[Proposition~3.15(c), Lemma~3.17 and
Proposition~3.18]{GonzalezWoodward}.  Its insertion therefore factors through
one of the restrictions in (8.2) and vanishes before division by the virtual
normal Euler class.  Only the two endpoints survive the virtual Kalkman
identity.  Their signs and node factors agree with Woodward's localized
adiabatic identity, giving (8.3).  Polarizing in the ordinary inputs gives
the entire row rather than only the scalar potential.
\end{proof}

For a birational map between smooth projective varieties, the required class
exists globally.  W{\l}odarczyk's smooth birational cobordism admits a smooth
projective equivariant completion whose source and sink quotients are the two
endpoints \cite[Proposition~2(B')]{Wlodarczyk}.  Choose \(p\) in the common
birational open supplied as a trivialized cylinder by the construction.
There the cobordism is a trivial orbit cylinder.  The
equivariant Poincar\'e dual of the closure of that orbit restricts to the two
point classes and misses every intermediate fixed component.  It is the
class \(a_p\) in (8.2).

There is a terminology point in the source.  W{\l}odarczyk calls the open
cobordism \emph{projective} when it is quasiprojective.  We use the resolved
proper equivariant completion constructed in the proof of
Proposition~2(B'); the resolution is an isomorphism over the source, sink,
and common-open cylinder.  Thus the extreme quotients and the class \(a_p\)
are unchanged.

\subsection{Neutral tails and the finite connection gate}

The two maps in (8.1) initially belong to opposite Novikov completions.  The
comparison needed for the point-row Boolean is smaller than a common field
for the complete source.  Fix a finite ample-energy and bulk Artin quotient
of the equivariant Novikov ring.  For a primitive affine gauge direction
\(\delta\), all ordinary curve degrees and fixed graph types are then finite;
only the integer coordinate along \(\delta\) can remain unbounded.  We first
control those unbounded tails.

Call the direction neutral when
\[
 c_1^{\Gm}(TW)\cdot\delta=0.                                 \tag{8.5}
\]
We compare formal monodromy after the Rees substitution
\[
 Q^d\longmapsto z^{c_1^{\Gm}(TW)\cdot d}Q^d.                 \tag{8.6}
\]

\begin{proposition}[Gamma-ratio reduction]
\label{prop:gamma-ratio-reduction}
At any finite Artin level, every individual rotation-fixed graph coefficient
in a neutral direction is a finite nilpotent derivative of a vector-valued
balanced Gamma-ratio kernel.  Its signed slope is
\(c_1^{\Gm}(TW)\cdot\delta\).  In a nonneutral direction, every homogeneous
coefficient is Laurent-finite.
\end{proposition}

\begin{proof}
Woodward's classification of the rotation-fixed graph locus is exhaustive
for sufficiently large area \cite[Corollary~9.10]{WoodwardQKIII}.  His
virtual-normal splitting gives
\[
 \operatorname{Eul}(N_\pm)=
 \operatorname{Eul}\!\left(
  (R\pi_*\operatorname{ev}^*T(W/\Gm))_{\mathrm{mov}}
 \right)(\mp\zeta)(\mp\zeta-\psi).                           \tag{8.7}
\]
The last two factors smooth the node and move the attaching point.  They are
independent of the affine gauge degree, so gluing neither removes one end of
the degree lattice nor creates a degree-dependent pole.  Formula (8.7)
applies when a nontrivial attached component is present; in the irreducible
unstable type those factors are absent and only the moving index remains
\cite[Example~9.15]{WoodwardQKIII}.

Pass to a finite cover which clears stabilizer denominators and apply the
splitting principle.  A virtual line in the moving index of the principal
weighted component has nilpotent Chern root \(\alpha_a\) and degree
\[
 n_a(k)=h_ak+s_a.                                             \tag{8.8}
\]
For every integer \(n\), its inverse index Euler factor is the single
meromorphic expression
\[
 \frac{1}{\prod_{m=0}^{n}(\alpha_a+m\zeta)}
 =\zeta^{-n-1}
   \frac{\Gamma(\alpha_a/\zeta)}
        {\Gamma(\alpha_a/\zeta+n+1)}.                         \tag{8.9}
\]
For \(n<0\), the left side denotes the Euler class of \(H^1\); for
\(n\geq0\), it denotes the inverse Euler class of \(H^0\).  Thus the two
degree rays are not assigned unrelated Gamma products.

The signed sum of the moving slopes is
\[
 \sum_a\epsilon_a h_a=c_1^{\Gm}(TW)\cdot\delta.              \tag{8.10}
\]
The gauge Lie algebra has slope zero.  Moving modes on attached fixed bubbles
have fixed rank at the chosen ordinary degree.  Their Euler factors are
finite products of linear functions of \(k\), possibly inverted.  Each
inverse factor is an adjacent Gamma ratio
\[
 (hk+a)^{-1}=\frac{\Gamma(hk+a)}{\Gamma(hk+a+1)},              \tag{8.11}
\]
so it has zero signed slope and does not alter (8.10).  Equivariant inputs
are polynomial in \(k\); Chern roots and psi classes are nilpotent at Artin
level; Liouville classes supply the Fourier exponential.

Under (8.5), equations (8.7)--(8.11) therefore give a balanced Gamma-ratio
kernel coefficientwise.  Polynomial equivariant insertions are Fourier
derivatives, while nilpotent Chern-root and bulk shifts give only finitely
many parameter derivatives.  For comparison, one moving-character residue
has the exact normalization
\[
 \operatorname*{Res}_{h\sigma+\alpha=-m}
 \Gamma(h\sigma+\alpha)\,d\sigma
 =\frac{(-1)^m}{h\,m!}.                                      \tag{8.12}
\]
The factorial is the moving-mode Euler factor, \(1/h\) is the
one-character Jacobian factor, and the sign is the complex orientation.  We
use (8.12) only as the normalization of a single simple pole.  Coincident or
higher poles require derivatives and the full localization expression,
including graph automorphism factors; no product-of-simple-residues formula
is asserted here.

If (8.5) fails, the virtual-dimension equation fixes \(k\) in each homogeneous
coefficient, hence the coefficient is Laurent-finite.  This proves the last
assertion.
\end{proof}

\begin{remark}[Exact analytic boundary]
\label{rem:neutral-boundary}
Aleshkin--Liu prove the relevant contour statement for a
finite-dimensional linear abelian GLSM with adjacent secondary-fan chambers,
a circuit, Calabi--Yau charges, and a grade-restriction window
\cite[Definition~5.18 and Theorem~5.21]{AleshkinLiu}.  Proposition
\ref{prop:gamma-ratio-reduction} does not construct those data for a nonlinear
virtual fixed graph.  Their theorem therefore does not determine the marked
connection matrices in Hypothesis~\ref{hyp:marked-threshold} below.

Gonz\'alez--Woodward's crepant wall-crossing theorem gives a weaker
distributional statement: neutral Picard shifts produce sums of derivatives
of delta distributions
\cite[Remark~1.18(d), Remark~4.6, and equation~(47)]{GonzalezWoodward},
and the authors explicitly do not deduce analytic continuation without an
independent convergence input \cite[Remark~4.7]{GonzalezWoodward}.  This does not replace
Hypothesis~\ref{hyp:marked-threshold}.  Proposition
\ref{prop:punctual-corner} exhibits the underlying obstruction: information
supported at a Fourier boundary can reappear with full support after inverse
transform, so almost-everywhere equality does not identify the marked formal
packet.
\end{remark}

\begin{proposition}[Clutching-tail holonomicity]
\label{prop:clutching-tail-holonomicity}
Fix a finite Artin level, an ordinary bubble degree, and a rotation-fixed
graph type.  After splitting into finitely many stabilizer congruence classes
and removing finitely many sign, stability, and zero-mode thresholds, every
point-marked endpoint tail is a finite nilpotent derivative of one
Gamma-ratio recurrence with affine arguments.  Its generating series is
holonomic.  In a neutral direction its coefficients have polynomial-logarithmic
growth after one exponential rescaling, so the tail has tempered radial
boundary values.
\end{proposition}

\begin{proof}
Corollary~9.10 and equations (54), (56), and (57) of
\cite{WoodwardQKIII} express a rotation-fixed component as a fibre product of
two clutching factors over the inertia quotient.  Each factor already
contains the complete ordinary stable-map space attached at that end.  Thus
arbitrary bubble trees occur in fixed proper stable-map factors once their
ordinary degrees and markings are fixed; no induction over bubble-tree edges
is required.

For \(G=\Gm\), an affine clutching cocharacter is an integer.  Away from
zero, its limit stratum depends only on its sign, while its inertia label
depends only on its residue modulo the finite stabilizer order.  A finite
sign/congruence split therefore fixes the coefficient stack and all
evaluation maps.  The stability type of the principal component is constant
on each resulting tail.  Stable tails retain the node and attaching factors
in (8.7); contracted tails use the unstable irreducible convention and
retain only the moving index.  The transition degrees form a finite set.

We also need the fixed obstruction theory, not only the fixed stack.  Apply
the normalization triangle to the pullback of the equivariant tangent
complex of \(W/\Gm\).  The clutching line is trivial on every attached
bubble, so the bubble complexes and node matching maps are independent of
the affine degree.  On the principal component, apply the splitting
principle.  A character root of weight \(h\) has degree \(hk+s\).  Roots
with \(h=0\) are constant.  For \(h\ne0\), the weights of its
\(H^0-H^1\) index form one consecutive affine interval.  Between consecutive
thresholds the zero-weight summand has constant rank and evaluation at a
generic point identifies it with the same root bundle on the coefficient
stack.  The \(H^1\) assertion is dual by Serre duality.  Any nonzero
clutching integer in the infinitesimal gauge map is removed by rescaling that
summand.  Hence the fixed normalization complex, its maps, its virtual class,
and its evaluations are constant on the tail.

The complementary moving weights give (8.9), with nilpotent Chern roots on
the fixed coefficient stack.  Only finitely many parameter derivatives are
therefore needed.  Gamma recurrence gives a rational finite-state recurrence
in \(k\), hence a holonomic generating series.  Under (8.5), Stirling's
formula cancels the \(k\log k\) terms in the signed product; the remaining
exponential is removed by one rescaling and the coefficients grow by a power
of \(k\) and \(\log k\).  This is exactly tempered Fourier growth.  Finite
threshold terms are Laurent polynomials and do not alter the conclusion.
\end{proof}

Tailwise holonomicity does not identify the two selected solutions.  The
bilateral sequence
\[
 c_k=\begin{cases}2^k,&k\geq0,\\3^{-k},&k<0\end{cases}
\]
has oriented germs
\[
 F_0(x)=\frac1{1-2x},\qquad F_\infty(x)=\frac3{x-3}.
\]
Both tails are first-order hypergeometric and tempered on their natural
boundary circles.  Away from the single threshold, the bilateral sequence
is annihilated by \((E-2)(E-1/3)\); multiplying by a polynomial in the index
absorbs the finite defect.  Nevertheless the two displayed rational
functions are not branches of one meromorphic function.  Thus neither a
zero wall term nor a polynomial recurrence determines the finite connection
data across a threshold.

\begin{hypothesis}[Marked threshold compatibility]
\label{hyp:marked-threshold}
At every finite Artin level and fixed ordinary degree, let \(C_j\) be the
connection matrix between two adjacent neutral clutching-tail systems.  Take
as state vector the pullback of the Gamma point row on finitely many common
equivariant inputs, together with its \(z\)-covariant derivatives until the
generated cyclic space stabilizes.  For every threshold,
\[
 r_jC_j=r_j,                                                  \tag{8.13}
\]
where \(r_j\) is this marked cyclic row.  These identifications commute with
input and bulk derivatives, quotient maps between Artin levels, and the Rees
substitution (8.6).
\end{hypothesis}

Condition (8.13) allows arbitrary connection on directions annihilated by
the point row.  A one-object Gamma/window comparison would imply it: a
window mutation is supported on the wall, whereas the skyscraper of a point
in the common open is unchanged and Euler-orthogonal to wall-supported
objects.  Pairing preservation and formal exponential labels alone do not
imply (8.13), as Sections~\ref{sec:incomplete-gamma}--\ref{sec:two-wall}
show.

\begin{lemma}[Cyclic-row spectral support]
\label{lem:cyclic-row-support}
Let \(T\) be an endomorphism of a finite dimensional vector space \(V\), and
let \(r\in V^*\).  Put
\[
 C_T(r)=\operatorname{span}\{r,rT,rT^2,\ldots\}\subset V^*.
\]
For every eigenvalue \(\lambda\), the restriction of \(r\) to the generalized
\(\lambda\)-eigenspace of \(T\) is nonzero if and only if the
\(\lambda\)-primary part of \(C_T(r)\) is nonzero.
\end{lemma}

\begin{proof}
Let \(e_\lambda(T)\) be the polynomial B\'ezout projector onto the generalized
\(\lambda\)-eigenspace.  The relevant primary projection of the row is
\(re_\lambda(T)\).  It belongs to \(C_T(r)\) and vanishes exactly when the
restriction of \(r\) does.
\end{proof}

\subsection{The finite packet}

Formal smooth-endpoint quantum Kirwan surjectivity is elementary in this
setting.  Modulo the positive-energy ideal, \(D\kappa_\pm\) is the derivative
of the classical Kirwan map and is surjective.  Choose a linear section of
the reduction.  Its arbitrary lift has composite \(1+N\), with \(N\) in the
augmentation ideal, and the convergent Neumann series
\(\sum_{j\geq0}(-N)^j\) gives a right inverse.  Since
\(D\tau_{Y_\pm,-}\) is an invertible fundamental solution, the maps (8.1)
are surjective as well.

The Rees factor in (8.6) is needed to compare formal monodromy.  Homogeneity
of \(D\kappa\) gives
\[
 \mu_YA-A\mu_W=\frac12A-\Theta_QA,                            \tag{8.14}
\]
because the acting group has dimension one and
\(\Theta_Q(Q^d)=c_1^{\Gm}(TW)\cdot d\,Q^d\).  After (8.6), the
\(\Theta_Q\)-term becomes \(z\partial_z\), so both endpoint maps intertwine
the \(z\)-connections with the same half-Tate shift.  The Euler-product term
intertwines because the virtual-dimension identity makes the quantum Kirwan
map homogeneous and its derivative a quantum-product morphism.  The target
bulk point \(\kappa(0)\) has positive energy and is formally parallel to the
large-radius point; this transport preserves the point row.

Choose finitely many common equivariant inputs whose images under each
surjection in (8.1) span the corresponding endpoint module; take the union
of two lift sets if necessary.  Evaluation on these images embeds the dual
endpoint module into one finite coordinate space.  After (8.6), this
embedding intertwines the dual \(z\)-connections up to the same half-Tate
shift.  The point row and its successive \(z\)-covariant derivatives
therefore generate a finite cyclic submodule in that coordinate space.

Proposition~\ref{prop:support-collapse} identifies the marked coefficient
functions before either endpoint completion is taken.
Proposition~\ref{prop:clutching-tail-holonomicity} supplies their holonomic
tail systems.  Hypothesis~\ref{hyp:marked-threshold} identifies the marked
solutions across every finite threshold, while nonneutral directions are
Laurent-finite by Proposition~\ref{prop:gamma-ratio-reduction}.  Thus the two
intrinsic endpoint cyclic row modules are conservative realizations of one
continued marked cyclic module.  No continuation of the complementary
equivariant quantum source is used.

\begin{theorem}[Conditional birational invariance of the point-row primary Boolean]
\label{thm:birational-point-primary}
Let \(Y_-\) and \(Y_+\) be smooth projective birational complex varieties.
Suppose a W{\l}odarczyk completion for this birational map satisfies
Hypothesis~\ref{hyp:marked-threshold} in every neutral affine gauge
direction.  Then, for every formal-monodromy eigenvalue \(\lambda\),
\[
 \rrow_{Y_-}|_{\cP_\lambda(Y_-)}\ne0
 \quad\Longleftrightarrow\quad
 \rrow_{Y_+}|_{\cP_\lambda(Y_+)}\ne0.                       \tag{8.15}
\]
Here the packets are realized through the marked cyclic continuation in the
hypothesis; no continuation of the complementary source is part of the
statement.
\end{theorem}

\begin{proof}
Use the global W{\l}odarczyk completion and its orbit-cylinder class.  By
the preceding construction, the two endpoint point rows generate the same
marked cyclic module under the dual formal-monodromy action.  Apply
Lemma~\ref{lem:cyclic-row-support}: the set of primary eigenvalues detected by
the marked row depends only on that cyclic module.  This gives the forward
implication, and the reverse implication is symmetric.
\end{proof}

Under Hypothesis~\ref{hyp:marked-threshold},
Theorem~\ref{thm:birational-point-primary} is the global replacement for
Hypothesis~\ref{hyp:rank-zero-target}.  Without that analytic input, the
two-tail example above and the earlier countermodels show why the conclusion
cannot be obtained from tailwise holonomicity, localized one-wall receivers,
or unmarked Fourier boundary data alone.
